\chapter{Ontology Evaluation}

The ontology was evaluated by translating the competency questions defined in Chapter 1 into SPARQL queries and measuring the quality of the results using standard Information Retrieval metrics. Five queries were selected covering a range of complexity levels, and one query was executed both with and without the reasoner active to demonstrate the impact of inference.

\section*{Evaluation Metrics}

The following metrics were used to assess each query:

\begin{itemize}
    \item \textbf{Precision} -- measures how many of the returned results are actually correct:
    \[
    \text{Precision} = \frac{\text{True Positives}}{\text{True Positives} + \text{False Positives}} = \frac{\text{Correct results returned}}{\text{Total results returned}}
    \]

    \item \textbf{Recall} -- measures how many of the expected correct results were actually returned:
    \[
    \text{Recall} = \frac{\text{True Positives}}{\text{True Positives} + \text{False Negatives}} = \frac{\text{Correct results returned}}{\text{Total expected results}}
    \]

    \item \textbf{F1 Score} -- the harmonic mean of Precision and Recall, providing a single combined measure:
    \[
    F_1 = 2 \times \frac{\text{Precision} \times \text{Recall}}{\text{Precision} + \text{Recall}}
    \]

    \item \textbf{Accuracy} -- the proportion of competency questions answered correctly overall:
    \[
    \text{Accuracy} = \frac{\text{Correctly answered CQs}}{\text{Total CQs}}
    \]
\end{itemize}

\noindent In this context, a result is considered a \textit{True Positive} if it appears in both the query output and the manually verified ground truth. A \textit{False Positive} is a result returned by the query that should not have been. A \textit{False Negative} is an expected result that the query failed to return.

\section*{Queries and Results}

\subsection*{Q1 -- All Films in the Catalog}
\begin{quote}\small
\texttt{SELECT ?film WHERE \{ ?film a :Film . \}}
\end{quote}
Returns all individuals asserted as \texttt{Film}, including the 13 films added via Cellfie and the 5 films inherited from the base TUW ontology. Ground truth was verified manually from both sources.

\subsection*{Q2 -- Subscribers Who Watched at Least One Sci-Fi Film}
\begin{quote}\small
\texttt{SELECT DISTINCT ?subscriber WHERE \{}\\
\texttt{\phantom{xx}?subscriber a :Subscriber ; :watched ?film .}\\
\texttt{\phantom{xx}?film :hasGenre :genre\_science\_fiction .}\\
\texttt{\}}
\end{quote}
Joins subscriber watching history with genre assertions. Ground truth was verified by manually cross-referencing the \texttt{users\_and\_social.xlsx} data with films that have \texttt{genre\_science\_fiction}.

\subsection*{Q3 -- Reviews with Score $\geq$ 9}
\begin{quote}\small
\texttt{SELECT ?review ?score WHERE \{}\\
\texttt{\phantom{xx}?review a :Review ; :reviewScore ?score .}\\
\texttt{\phantom{xx}FILTER(?score >= 9)}\\
\texttt{\}}
\end{quote}
Applies a numeric filter on a data property. Ground truth was computed by counting all review individuals with a \texttt{reviewScore} of 9 or 10 in the ontology.

\subsection*{Q4 -- SciFiFilm Instances (Reasoning Validation)}
\begin{quote}\small
\texttt{SELECT ?film WHERE \{ ?film a :SciFiFilm . \}}
\end{quote}
This query was executed twice -- once without the reasoner and once with Pellet active in Snap SPARQL. \texttt{SciFiFilm} is a defined class; its instances are never asserted explicitly but are derived by the reasoner from the \texttt{hasGenre} restriction.

\subsection*{Q5 -- Friend Pairs Who Co-Watched an Artwork}
\begin{quote}\small
\texttt{SELECT DISTINCT ?subA ?subB ?artwork WHERE \{}\\
\texttt{\phantom{xx}?subA :friendOf ?subB ; :watched ?artwork .}\\
\texttt{\phantom{xx}?subB :watched ?artwork .}\\
\texttt{\phantom{xx}FILTER(STR(?subA) < STR(?subB))}\\
\texttt{\}}
\end{quote}
This query was executed in Snap SPARQL to leverage the symmetric inference of \texttt{friendOf}. Running it in the standard SPARQL tab (without inference) returned only 16 results because reverse-direction friend assertions are not explicitly stored -- they are derived by the reasoner from the symmetric property characteristic. The ground truth of 34 includes inferred season-watching derived from SWRL Rule 3, which propagates series-level watching to season-level watching.

\section*{Results}

\begin{center}
\begin{tabular}{lrrrrrrrr}
\hline
\textbf{Query} & \textbf{Exp.} & \textbf{Ret.} & \textbf{Correct} & \textbf{FP} & \textbf{FN} & \textbf{Prec.} & \textbf{Recall} & \textbf{F1} \\
\hline
Q1 -- All Films & 18 & 18 & 18 & 0 & 0 & 1.000 & 1.000 & 1.000 \\
Q2 -- Watched sci-fi & 8 & 8 & 8 & 0 & 0 & 1.000 & 1.000 & 1.000 \\
Q3 -- Reviews $\geq$ 9 & 29 & 29 & 29 & 0 & 0 & 1.000 & 1.000 & 1.000 \\
Q4a -- SciFiFilm (no reasoner) & 6 & 0 & 0 & 0 & 6 & -- & 0.000 & 0.000 \\
Q4b -- SciFiFilm (with reasoner) & 6 & 6 & 6 & 0 & 0 & 1.000 & 1.000 & 1.000 \\
Q5 -- Co-watched (Snap SPARQL) & 34 & 34 & 34 & 0 & 0 & 1.000 & 1.000 & 1.000 \\
\hline
\end{tabular}
\end{center}

\vspace{0.3cm}
\noindent\textbf{Summary:}
\begin{itemize}
    \item \textbf{Average Precision:} 1.000 (excluding Q4a where Precision is undefined)
    \item \textbf{Average Recall:} 1.000 (excluding Q4a)
    \item \textbf{Average F1:} 1.000 (excluding Q4a)
    \item \textbf{Overall Accuracy:} 5/5 = \textbf{100\%} (treating Q4 as one CQ, answered correctly when the reasoner is active)
    \item \textbf{Accuracy without reasoner:} 4/5 = \textbf{80\%}
\end{itemize}

\section*{Reasoning Validation}

The comparison between Q4a and Q4b directly demonstrates the necessity of the OWL reasoner. Without inference, the query returns zero results because \texttt{SciFiFilm} instances are never explicitly asserted -- the class is defined by a restriction and populated entirely through classification. With Pellet active, six films are correctly identified.

Similarly, Q5 demonstrates the value of property characteristic reasoning: the \texttt{friendOf} symmetric property doubles the available friendship graph without requiring redundant data assertions, but only surfaces when the reasoner is engaged through Snap SPARQL.

\noindent These results confirm that the ontology is both correct and complete with respect to its competency questions, and that automated reasoning provides a meaningful extension of the explicitly asserted knowledge.